% !TeX root = practica_semanas_01_03.tex
% !TeX program = pdflatex
% Práctica acumulativa de las semanas 01 a 03.
\documentclass[10pt,a5paper]{article}
\newif\ifmostrarSoluciones
\mostrarSolucionestrue
\usepackage{estilo_apuntes}
\encabezado{Semana 4}{Práctica acumulativa}

\begin{document}
\tituloclase{Práctica de las semanas 1 a 3}{Funciones, composición y límites}{Semana 4 · 1 hora}

\begin{objetivos}
\begin{itemize}
  \item Integrar el análisis de dominios, composiciones, transformaciones y biyectividad.
  \item Utilizar límites laterales, simplificación algebraica, racionalización y el teorema del Sanduche.
  \item Justificar cada respuesta indicando las restricciones y propiedades empleadas.
\end{itemize}
\end{objetivos}

\begin{nota}[organización sugerida]
Trabaje aproximadamente 12 minutos en cada ejercicio y reserve los últimos 12 minutos para comparar estrategias y corregir los pasos principales. No basta escribir una respuesta final: deben justificarse los dominios, las simplificaciones y las conclusiones sobre límites.
\end{nota}

\begin{ejercicios}[práctica acumulativa]
\begin{enumerate}
  \item Sean
  \[
    f\colon(-\infty,4]\longrightarrow[0,\infty),
    \qquad
    f(x)=\sqrt{4-x},
  \]
  y
  \[
    g\colon\mathbb{R}\setminus\{1\}\longrightarrow\mathbb{R},
    \qquad
    g(x)=\frac{1}{x-1}.
  \]
  Halle las leyes de asignación de $g\circ f$ y $f\circ g$, y determine cuidadosamente el dominio de cada composición.

  \item Considere la función
  \[
    h\colon[1,\infty)\longrightarrow[-3,\infty),
    \qquad
    h(x)=2(x-1)^2-3.
  \]
  Describa las transformaciones que permiten obtener su gráfica a partir de $y=x^2$. Analice si $h$ es biyectiva.

  \item Sea
  \[
    F\colon[-3,1)\cup(1,\infty)\longrightarrow\mathbb{R}
  \]
  la función definida por
  \[
    F(x)=
    \begin{cases}
      \dfrac{\sqrt{x+3}-2}{x-1}, & x<1,\\[4mm]
      \dfrac{x^2+x-2}{x-1}, & x>1.
    \end{cases}
  \]
  Decida si existe $\displaystyle\lim_{x\to1}F(x)$.

  \item Calcule, justificando cada simplificación,
  \[
    \lim_{x\to4}
    \frac{\sqrt{x+5}-3}{x^2-16}.
  \]
  Indique por qué no es válido sustituir directamente.
\end{enumerate}
\end{ejercicios}

\ifmostrarSoluciones
\begin{solucion}[de la práctica acumulativa]
\begin{enumerate}
  \item La primera composición tiene ley
  \[
    (g\circ f)(x)
    =\frac{1}{\sqrt{4-x}-1}.
  \]
  Se requiere $x\leq4$ para evaluar $f(x)$ y $f(x)\neq1$ para poder aplicar $g$. Como
  \[
    \sqrt{4-x}=1
    \iff x=3,
  \]
  resulta
  \[
    D_{g\circ f}=(-\infty,4]\setminus\{3\}.
  \]
  Para la segunda composición,
  \[
    (f\circ g)(x)
    =\sqrt{4-\frac{1}{x-1}}.
  \]
  Se requiere $x\neq1$ y
  \[
    4-\frac{1}{x-1}
    =\frac{4x-5}{x-1}
    \geq0.
  \]
  El análisis de signos en los puntos $1$ y $5/4$ da
  \[
    D_{f\circ g}
    =(-\infty,1)\cup\left[\frac54,\infty\right).
  \]

  \item A partir de $y=x^2$, la gráfica se traslada una unidad a la derecha, se estira verticalmente por un factor de $2$ y se traslada tres unidades hacia abajo. En el dominio $[1,\infty)$, la función es estrictamente creciente y su imagen es $[-3,\infty)$; por tanto, es inyectiva y sobreyectiva sobre el codominio declarado.

  \item Para $x<1$ y $x\neq1$, racionalizamos:
  \[
    \frac{\sqrt{x+3}-2}{x-1}
    =\frac{1}{\sqrt{x+3}+2}.
  \]
  Por tanto,
  \[
    \lim_{x\to1^-}F(x)=\frac14.
  \]
  Para $x>1$, factorizamos:
  \[
    \frac{x^2+x-2}{x-1}
    =\frac{(x-1)(x+2)}{x-1}
    =x+2.
  \]
  En consecuencia,
  \[
    \lim_{x\to1^+}F(x)=3.
  \]
  Como los límites laterales son diferentes, el límite bilateral no existe.

  \item La sustitución directa produce la forma indeterminada $0/0$. Para $x\neq4$, racionalizamos y factorizamos:
  \begin{align*}
    \frac{\sqrt{x+5}-3}{x^2-16}
    &=\frac{x-4}
      {(\sqrt{x+5}+3)(x^2-16)}\\
    &=\frac{x-4}
      {(\sqrt{x+5}+3)(x-4)(x+4)}\\
    &=\frac{1}
      {(\sqrt{x+5}+3)(x+4)}.
  \end{align*}
  La última expresión es localmente igual a la original alrededor de $4$, excepto en $x=4$. Por sustitución,
  \[
    \lim_{x\to4}
    \frac{\sqrt{x+5}-3}{x^2-16}
    =\frac{1}{(3+3)(4+4)}
    =\frac{1}{48}.
  \]

\end{enumerate}
\end{solucion}
\fi

\end{document}
